% !TeX document-id = {1009e7ce-5d51-4851-b71b-515f3f53617f}
% !TeX TXS-program:compile = txs:///pdflatex/[--shell-escape]
\documentclass[spanish,addpoints,answers,a4paper]{exam}
\usepackage[T1]{fontenc}
\usepackage[spanish]{babel}
\usepackage{libertine}
\usepackage[tmargin=2cm,rmargin=2cm,lmargin=2cm]{geometry}
\usepackage{amsmath,bm}
\usepackage[shortlabels]{enumitem}
\usepackage{siunitx}
\sisetup{
	group-four-digits = true,
	group-separator = {,}
}
\usepackage{multicol}

\usepackage{graphicx}
\graphicspath{{../../../images/}}
\pagestyle{plain}

\usepackage{minted}
\usepackage{hyperref}
\renewcommand{\solutiontitle}{\noindent\textbf{Solución}\par\noindent}
\renewcommand\listingscaption{Listado}
\newcommand{\unmarkedfntext}[1]{%
	\begingroup
	\renewcommand\thefootnote{}\footnote{#1}%
	\addtocounter{footnote}{-1}%
	\endgroup
}
\spanishdatedel
\decimalpoint

\def\LOGOUNI{%
	\begin{picture}(0,0)\unitlength=1cm
	\put (-13.5,-2.5) {\includegraphics[width=2.8cm]{logouni}}
	\end{picture}
}

\def\LOGOCTIC{%
	\begin{picture}(0,0)\unitlength=1cm
	\put (-8,-1.2) {\includegraphics[height=1cm]{logocticblack}}
	\end{picture}
}

\newcommand{\mychar}{%
	\begingroup\normalfont
	\includegraphics[height=\fontcharht\font`\B]{Octocat.png}%
	\endgroup
}

\begin{document}

\begin{center}
	\sffamily\bfseries\scshape
	{\Large UNIVERSIDAD NACIONAL DE INGENIERÍA}\LOGOUNI\\
	Centro de Tecnologías de la Información y Comunicaciones\LOGOCTIC\\
\end{center}

\vspace{.8cm}

\begin{center}\sffamily\bfseries\large
	Programación en \texttt{C++} \\
	TEMA $\bm{1}$: Entrada y salida de datos
\end{center}

\vspace{.5cm}
\noindent
\makebox[\textwidth]{Nombres y apellidos:\enspace\hrulefill}
\vspace{0.2in}
\makebox[\textwidth]{Nombres y apellidos del instructor:\enspace MSc. Víctor Melchor Espinoza.\hfill}

\section*{Ejercicios de práctica}

\begin{questions}

\question Elabore un programa que envíe a pantalla el saludo ``\textbf{Hola mundo}''.

\begin{solution}
	\begin{listing}[H]
		\footnotesize
		\inputminted{cpp}{1.cc}
		\caption{Programa \texttt{1.cc}.}
		\label{lst:1.1}
	\end{listing}
\end{solution}

\question Presente un ejemplo del uso de los operadores de asignación y aritméticos.

\begin{solution}
	\begin{listing}[H]
		\footnotesize
		\inputminted{cpp}{2.cc}
		\caption{Programa \texttt{2.cc}.}
		\label{lst:1.2}
	\end{listing}
\end{solution}

\question Determinar el salario neto de un trabajador que trabajó $36$ horas, con un pago por hora de \SI{50.00}[S/]{} soles y la tasa del impuesto del $13$\%.\label{q:3}

\begin{solution}
	\begin{listing}[H]
		\footnotesize
		\inputminted{cpp}{3.cc}
		\caption{Programa \texttt{3.cc}.}
		\label{lst:1.3}
	\end{listing}
\end{solution}

\question Resolver el ejercicio~\ref{q:3} para el caso general, conociendo el número de horas trabajadas, pago por hora y la tasa de impuesto.

\begin{solution}
	\begin{listing}[H]
		\footnotesize
		\inputminted{cpp}{4.cc}
		\caption{Programa \texttt{4.cc}.}
		\label{lst:1.4}
	\end{listing}
\end{solution}

\question Escriba un programa que pida dos números enteros que representen al dividendo y divisor. Obtenga el cociente y el resto.

\begin{solution}
	\begin{listing}[H]
		\footnotesize
		\inputminted{cpp}{5.cc}
		\caption{Programa \texttt{5.cc}.}
		\label{lst:1.5}
	\end{listing}
\end{solution}

\question Dado un número de tres cifras, hallar la suma de sus cifras.

\begin{solution}
	\begin{listing}[H]
		\footnotesize
		\inputminted{cpp}{6.cc}
		\caption{Programa \texttt{6.cc}.}
		\label{lst:1.6}
	\end{listing}
\end{solution}

\question Pida un valor del radio y obtenga el área del círculo.

\begin{solution}
		\begin{listing}[H]
		\footnotesize
		\inputminted{cpp}{7.cc}
		\caption{Programa \texttt{7.cc}.}
		\label{lst:1.7}
	\end{listing}
\end{solution}

\question Elabore un programa que pida un monto en dólares y obtenga el monto equivalente en soles. Defina el tipo de cambio como una constante.

\begin{solution}
	\begin{listing}[H]
		\footnotesize
		\inputminted{cpp}{8.cc}
		\caption{Programa \texttt{8.cc}.}
		\label{lst:1.8}
	\end{listing}
\end{solution}

\end{questions}

\section*{Ejercicios obligatorios}

\begin{questions}

\question En el curso de \textbf{Programación}, el promedio final se compone de los siguientes porcentajes:

\begin{itemize}[--]
	\item $60$\% del promedio de cuatro exámenes (PE).
	\item $15$\% del promedio de laboratorios (PL)
	\item $15$\% del promedio de tareas (PT)
	\item $10$\% de la nota del proyecto final (PR)
\end{itemize}

Desarrolle un programa en el que solicite las notas de los cuatro exámenes, del promedio de laboratorios, del promedio de tareas y la nota del proyecto. En base a esto, el programa determinará la nota final del alumno.

\textbf{Ejemplo de salida:}
\begin{minted}{bash}
Bienvenido al programa que calcula la Nota Final del curso de Programación.
Ingrese la nota del Examen N°1: 10
Ingrese la nota del Examen N°2: 08
Ingrese la nota del Examen N°3: 09
Ingrese la nota del Examen N°4: 09
Ingrese el promedio de Laboratorios: 10
Ingrese el promedio de Tareas: 14
Ingrese la nota del Proyecto Final: 14

Reporte de notas:
Nota del Examen N°1: 10
Nota del Examen N°2: 08
Nota del Examen N°3: 09
Nota del Examen N°4: 09
Promedio de Exámenes (PE): 09
Promedio de Laboratorios (PL): 14
Promedio de Tareas (PT): 14
Nota del Proyecto Final (PF): 14
La fórmula de la Nota Final es: NF=0.60PE+0.15P.L+0.15PT+0.10PR
Vuestra Nota Final es: 11
\end{minted}

\begin{solution}
	\begin{listing}[H]
		\footnotesize
		\inputminted{cpp}{1.cpp}
		\caption{Programa \texttt{1.cpp}.}
		\label{lst:1.9}
\end{listing}
\end{solution}

\question Realiza un programa que proporcione el desglose en billetes y monedas de una cantidad entera de euros. Recuerda que hay billetes de $500$, $200$, $100$, $50$, $20$, $10$ y $5$ euros, también hay monedas de $2$ y $1$ euro.

\begin{minted}{bash}
Ingrese el importe total (en euros) = 977

1 billete de 500 euros
2 billetes de 200 euros
1 billete de 50 euros
1 billete de 20 euros
0 billete de 10 euros
1 moneda de 5 euros
2 monedas de 1 euro
\end{minted}

\begin{solution}
	\begin{listing}[H]
		\footnotesize
		\inputminted{cpp}{2.cpp}
		\caption{Programa \texttt{2.cpp}.}
		\label{lst:1.10}
\end{listing}
\end{solution}

\end{questions}

\begin{flushright}\bfseries
Centro de Tecnologías de la Información y Comunicaciones (CTIC)\\[2mm]
\today\unmarkedfntext{Código disponible en \href{https://github.com/carlosal1015/Cpp-Programming}{\mychar{}}.}
\end{flushright}

\end{document}